\section{שלוש שכבות של סוכן: Software, Skill, Agent}

אחת התרומות המרכזיות של ההרצאה היא מודל השכבות השלוש, המסביר כיצד סוכן
\eng{AI} בנוי מרכיבים שונים בעלי אחריות ברורה. המודל מאפשר להבין היכן
להציב לוגיקה קבועה, היכן להשתמש בשפה טבעית ואיפה ``המוח'' של הסוכן פועל.

% Three-layer agent architecture
\begin{figure}[H]
\centering
\begin{tikzpicture}[
  layer/.style={draw=primary, fill=primary!8, rounded corners=4pt,
    minimum width=9cm, minimum height=1.2cm, align=center, font=\small}
]
  \node[layer, fill=accentgreen!15, draw=accentgreen] (agent) at (0,2.4)
    {\textbf{שכבת Agent}\\\eng{LLM + Memory + Context + RAG + Tools + Skills}};
  \node[layer, fill=accentblue!12, draw=accentblue] (skill) at (0,0.8)
    {\textbf{שכבת Skill}\\עטיפת שפה טבעית מעל קוד תבניתי};
  \node[layer, fill=lightgray, draw=secondary] (soft) at (0,-0.8)
    {\textbf{שכבת Software}\\פונקציות ותהליכים קבועים};
  \draw[-{Stealth[length=2.5mm]}, thick, secondary] (soft.north) -- (skill.south);
  \draw[-{Stealth[length=2.5mm]}, thick, secondary] (skill.north) -- (agent.south);
\end{tikzpicture}
\caption{שלוש שכבות הסוכן לפי חומר ההרצאה: תוכנה, סקיל וסוכן}
\label{fig:three-layers}
\end{figure}


\subsection{שכבת Software --- התשתית הקבועה}

בשכבה התחתונה נמצאת התוכנה התבניתית: פונקציות, אלגוריתמים ותהליכים שחייבים
לרוץ באופן דטרמיניסטי. דוגמאות: חישוב מרחק קוסינוס בין וקטורים, קריאת קובץ
\eng{CSV}, שליחת בקשת \eng{HTTP} או קומפילציה של \LaTeX{} ל-\eng{PDF}.
שכבה זו אינה ``מבינה'' שפה --- היא מבצעת הוראות מדויקות.

\subsection{שכבת Skill --- שפה טבעית מעל קוד}

שכבת ה-\eng{Skill} עוטפת את התוכנה במעטפת של שפה טבעית. במקום לספק
פרמטרים מדויקים לפונקציה, המשתמש או הסוכן מתארים בכוונה כללית מה נדרש,
וה-\eng{LLM} בשיתוף ה-\eng{Skill} ממפה את התיאור לפרמטרים הנכונים.
לדוגמה: ``חפש דירות זולות ברחוב רוטשילד'' במקום קואורדינטות מדויקות של
שדות בטופס.

\begin{insightbox}
\eng{Skill} הוא לא ``סוכן'' --- הוא מנגנון שמאפשר לדבר עם תוכנה קיימת
בשפה חופשית. הסוכן משתמש ב-\eng{Skills} כאשר הוא מזהה שמשימה מתאימה
לתיאור של אחד מהם.
\end{insightbox}

\subsection{שכבת Agent --- המוח וההחלטות}

שכבת הסוכן כוללת את ה-\eng{LLM}, את חלון ההקשר (\eng{Context Window}),
את הזיכרון, את ה-\eng{RAG}, את הכלים (\eng{Tools}) ואת היכולת לבחור
\eng{Skills} מתאימים. הסוכן מקבל מטרה, מתכנן, מפעיל כלים ומסכם תוצאות.
בניגוד ל-\eng{Skill} בודד, הסוכן יכול לשרשר מספר פעולות ולהתאים את עצמו
למשוב.

\agenttable{
\caption{השוואת שלוש השכבות}
\label{tab:three-layers}
\begin{tabular}{@{}p{2.2cm}p{3.5cm}p{3.5cm}p{3.5cm}@{}}
\toprule
\textbf{מאפיין} & \textbf{Software} & \textbf{Skill} & \textbf{Agent} \\
\midrule
שפה טבעית & לא & כן (ממשק) & כן (ליבה) \\
קביעות & גבוהה & בינונית & נמוכה \\
תכנון רב-שלבי & לא & מוגבל & כן \\
דוגמה & \eng{cosine()} & תרגום מסמך & חקירת אירוע \\
\bottomrule
\end{tabular}
}

\subsection{זרימת עבודה אופיינית}

כאשר משתמש שואל ``מה בירת מדינת ישראל?'', הסוכן:
\begin{enumerate}[rightmargin=2em]
  \item מנתח את הפרומפט ב-\eng{LLM}.
  \item מחפש \eng{Skill} או כלי חיפוש רלוונטי.
  \item מפעיל את שכבת ה-\eng{Software} (למשל חיפוש ב-\eng{RAG}).
  \item מרכיב תשובה מתוך חלון ההקשר ומחזיר למשתמש.
\end{enumerate}

\begin{warningbox}
בלבול בין שכבות --- למשל הצבת לוגיקה עסקית קריטית רק בפרומפט ללא
קוד תבניתי --- עלול לגרום לתוצאות לא עקביות. כלל אצבע: מה שחייב
להיות זהה בכל הרצה שייך לשכבת \eng{Software} או ל-\eng{CLS}.
\end{warningbox}

\subsection{יישום ארגוני}

בארגון, ניתן לייצר ספריית \eng{Skills} לכל מחלקה: כספים, משאבי אנוש,
\eng{CTI}. הסוכן המרכזי בוחר את ה-\eng{Skill} המתאים לפי תיאור בראש
הקובץ. כך נשמרת אחידות בלי לשכפל קוד בכל סוכן.
